\section{对 LLM Agent 的系统启发}

前面案例覆盖机器人、陌生队友、人类反馈和赛车系统；本节把共同结构迁移到 LLM Agent。迁移重点是可调试闭环、队友分布、反馈组合和行为治理，而不是把物理控制算法逐字搬到文本模型。

把 Stone 这节课映射回 2025 年 LLM Agent，至少有四个直接结论。

\begin{importantbox}{启发 1：Agent 评测必须包含 interaction outcome}
仅看单轮回答质量会误导。应评估：
\begin{itemize}
    \item 长链任务完成率；
    \item 与人/工具/其他 agent 的协同稳定性；
    \item 出错恢复路径是否可控。
\end{itemize}
\end{importantbox}

\begin{importantbox}{启发 2：分层动作空间可以显著降本}
Slack 的 latent action 思想在软件 Agent 里也成立：不要总是直接生成低层动作序列，可以先学高层 action primitives，再做下游任务优化。
\end{importantbox}

\begin{importantbox}{启发 3：数据策划和 replay 配比决定战术能力}
GT Sophy 证明了 ``看到什么分布，就学到什么策略''。在 coding/ops agent 里，同样需要精心构造困难场景、对抗场景和边界场景的训练比例。
\end{importantbox}

\begin{warningbox}{启发 4：安全约束应显式进入训练目标}
无论是机器人避障还是竞速 etiquette，本质都是 ``把合规从后处理挪到训练过程''。LLM Agent 也应把权限、审计、风险动作约束前置。
\end{warningbox}

\subsection{本章小结}
这节课不是在和 LLM Agent 竞争叙事，而是在提供一套更稳的工程母框架：定义闭环、分层建模、场景化训练、系统级评测。

